\chapter{Conclusion}
\label{chap:conclusion}

\par Doclane started from a straightforward observation: the process of submitting documents to a public institution is, in most cases, slower and more frustrating than it needs to be. Citizens take time off work to stand in queues, sometimes only to find out that one document is missing or incorrectly formatted, after which the whole cycle repeats. This thesis set out to address that by building a system that moves the document exchange entirely online, structures it clearly for both sides, and keeps the citizen informed throughout.

\par The result is a web-based platform where citizens can browse request templates, upload the required documents in designated slots, and track the status of their submission without needing to visit the institution. On the other side, department employees can claim incoming requests, review each document individually, and push feedback back to the citizen directly through the system if something needs to be corrected. The loop continues until all documents are approved, at which point the request is marked complete. Role-based access ensures that each request has a clear owner within the department and that sensitive information is only visible to those who are meant to see it.

\par The technical implementation reflects an effort to build something production-grade rather than just functional. The backend is written in Go, organized into clearly separated layers, and deployed as a containerized workload on EKS with autoscaling and health checks configured. The frontend is a Next.js application using server-side rendering to keep the experience fast despite the data-heavy nature of the pages. Infrastructure is defined entirely in Terraform and designed to be torn down and recreated on demand, which kept development costs low while maintaining a realistic deployment environment. AWS managed services handle authentication through Cognito, file storage through S3, and document processing through Textract, Bedrock, and Polly, with all credentials managed through short-lived IRSA roles rather than static keys.

\par The AI-assisted features sit alongside the core workflow rather than inside it. Textract extracts text from uploaded scans, Bedrock assesses whether a document matches what the template requires, and Polly can convert document content to audio. None of these produce decisions; they produce context that the reviewing employee can use or ignore. This boundary was drawn deliberately, both because the accuracy of these tools does not yet justify giving them decision authority, and because the legal responsibility for approving or rejecting a submission must remain with a person.

\par The system has limitations that are worth acknowledging plainly. It does not integrate with national identity infrastructure, so submissions carry no stronger identity verification than a self-registered email account. The dependence on AWS means that certain features are tied to a single provider in ways that would take effort to unwind. And as with any tool introduced into an institutional setting, the gap between a working system and a widely adopted one is largely organisational rather than technical. These are the directions in which the system would need to grow to move from prototype to production deployment at scale.

\section{Future improvements}

\par The two most meaningful extensions to the current system would address identity verification and in-person scheduling, both of which represent the remaining points of friction in an otherwise digital workflow.

\par Integration with eIDAS-compliant digital signing would be the higher-impact of the two. eIDAS is the EU regulatory framework for electronic identification and trust services, and compliance with it would allow users to verify their identity in a legally recognised way without appearing in person. This would remove the last remaining reason a citizen might need to visit the institution before their request is resolved, since document submission, review, and identity confirmation could all happen through the platform. It would also increase the legal weight of submissions, which is a prerequisite for the system to be used for higher-stakes administrative procedures.

\par The second improvement concerns what happens after all documents are approved. At that point, many types of requests still require the citizen to collect a physical document from the institution. Currently, coordinating that appointment happens outside the system entirely. A scheduling feature would allow the employee to propose available time slots once a request is complete, and the citizen would select one directly through the platform. This is a relatively contained feature to implement but it would close the last gap in the workflow and eliminate another category of phone calls and in-person coordination that institutions currently handle manually.

\par Together, these two additions would bring Doclane considerably closer to the goal it was built around: a document workflow that is fully digital from first submission to final collection, where physical presence is reduced to only what is legally or practically unavoidable.